\section{Introduction}

Connected autonomous vehicles (CAVs) benefit from cooperative perception:
sharing sensor data or intermediate feature representations across a
vehicular network transcends physical limitations such as occlusion,
limited sensor range, and adverse weather.  Recent cooperative
perception-prediction stacks such as CMP~\cite{cmp2024,cmp_ieee_2023} demonstrate that
aggregating multi-agent data yields measurably better tracking and
forecasting than single-agent systems.

Deploying these stacks at dense scale exposes two infrastructure
bottlenecks.  First, direct vehicle-to-vehicle (V2V) exchange of
intermediate feature maps is physically infeasible under dense
participation: our ns-3 5G-LENA NR~V2X sidelink simulations show a
Packet Reception Ratio (PRR) of zero for 200\,KB payloads at every
tested $N$, because a single intermediate feature map exceeds the
sidelink transport block capacity.  Second, while centralizing on a
Multi-Access Edge Computing (MEC) server resolves the bandwidth
constraint, it introduces a \emph{deadline cliff}: the full edge
pipeline (spatial-attention fusion, multi-object tracking, trajectory
prediction) scales super-linearly with the number of agents, crossing
a 100\,ms planner deadline at $N \approx 10$ on commodity hardware and
reaching 184\,ms at $N=64$.

The edge server possesses three assets unavailable to individual
vehicles: (i)~persistent Kalman-filtered track state across ticks,
enabling divergence-gated prediction reuse; (ii)~a global scene graph
containing all ego poses, planned paths, and fields of view; and
(iii)~a single centralized compute budget governed by one deadline.
These assets make the edge the natural locus for runtime adaptation.

We present a deadline-aware adaptive controller with two decision
variables: $K$, the number of agents retained after an input
contribution filter, and $x_i \in \{0,1\}$, a per-object binary
indicating whether tracked object $i$ receives fresh transformer
prediction or reuses a cached trajectory.  The controller implements
three mechanisms:
%
\begin{enumerate}
  \item \textbf{Input contribution filtering.} Score each agent by
    detection confidence within the ego-vehicle conflict zone; retain
    the top-$K$ agents that maximize estimated downstream utility.
  \item \textbf{Temporal amortization.} Gate fresh prediction inference
    on per-object divergence between the observed track state and the
    cached prediction, using velocity-normalized positional and heading
    thresholds.
  \item \textbf{Risk-budgeted scheduling.} Allocate the residual
    prediction budget (after fusion and tracking) to divergent objects
    ranked by a continuous risk score integrating speed, proximity to
    ego, occlusion status, and conflict geometry.
\end{enumerate}
%
Rate adaptation (reducing the global tick frequency) is retained as a
fallback triggered only when no feasible $K$ keeps the estimated
pipeline cost within the deadline.

We evaluate the controller offline on real post-backbone features
from the OPV2V dataset, using a WorldFusion (Where2Comm attention)
fusion model, AB3DMOT tracking, and MTR transformer prediction.
Across $N \in \{4,\ldots,64\}$, the adaptive controller maintains
100\% deadline compliance at a 100\,ms budget while the static baseline
fails at $N \geq 16$.  A fixed-ground-truth evaluation with explicit
miss-rate penalties and risk stratification confirms that prediction
quality for high-risk objects is preserved under adaptation.

Our contributions:
\begin{enumerate}
  \item \textbf{Dense-scale characterization.} Per-stage latency
    measurements for a WorldFusion\,+\,AB3DMOT\,+\,MTR edge pipeline,
    identifying the deadline cliff and its per-component drivers.
  \item \textbf{Adaptive controller.} A two-variable greedy controller
    (input filter $K$, per-object prediction $x_i$) that maintains
    deadline compliance through $N=64$ while preserving safety-critical
    prediction quality.
  \item \textbf{Risk-stratified evaluation methodology.} A fixed-GT
    evaluation protocol that penalizes missing predictions and
    separates quality reporting by object risk level.
\end{enumerate}
